\documentclass{article}
\usepackage{amsmath}
\title{Mathematical Operators}
\author{ }
\date{}

\begin{document}
\maketitle

\section*{1. Summation}
A summation adds a sequence of values.

\textbf{LaTeX source code:}
\begin{verbatim}
\[
\sum_{i=1}^{n} i
\]
\end{verbatim}
\textbf{Expected compiled output:}
\[
\sum_{i=1}^{n} i
\]
\textbf{Explanation:} \verb|\sum| creates the summation symbol. The subscript gives the starting value and the superscript gives the ending value.

\section*{2. Product}
A product multiplies a sequence of values.

\textbf{LaTeX source code:}
\begin{verbatim}
\[
\prod_{i=1}^{n} i
\]
\end{verbatim}
\textbf{Expected compiled output:}
\[
\prod_{i=1}^{n} i
\]

\section*{3. Limit}
A limit describes the value an expression approaches.

\textbf{LaTeX source code:}
\begin{verbatim}
\[
\lim_{x \to 0} f(x)
\]
\end{verbatim}
\textbf{Expected compiled output:}
\[
\lim_{x \to 0} f(x)
\]
\textbf{Explanation:} \verb|\lim| creates the limit operator and \verb|\to| creates the arrow.

\section*{4. Practice}
Write a summation from \(i=1\) to \(5\), and a limit as \(x\) approaches \(2\).
\end{document}
